\section{Operational lessons from deployment}

The operational deployment of ECHOREPO revealed several requirements that were less apparent when the infrastructure was initially designed. These lessons concern not only the implementation of ECHOREPO, but more generally the management of citizen-generated research data once they move beyond the collection stage.

First, data quality should not be treated exclusively as a filtering problem. In a citizen-science setting, an erroneous field does not necessarily invalidate the complete observation. A sample with an incorrect coordinate may still contain valid images, field observations and laboratory results. Preserving source identifiers and provenance makes it possible to investigate such cases and, where sufficient evidence exists, recover information rather than simply discard the record. Quality assurance should therefore support detection, classification and recovery as distinct operations.

Second, the participant-facing acquisition system and the long-term research repository have different responsibilities and should not necessarily be implemented as a single system. The acquisition application is optimised for interaction with participants, data entry and field operation, whereas the research repository must support provenance, harmonisation, enrichment, controlled access and long-term publication. Separating these functions allows the collection application to evolve without requiring corresponding changes to the public research-data model and reduces dependence on a particular mobile backend or authentication technology.

Third, the data associated with a citizen-science observation may change substantially during its lifecycle. A record that initially consists of a small set of field observations can later become associated with photographs, laboratory measurements and large molecular biodiversity datasets. In the ECHOREPO deployment, biodiversity data account for most of the overall storage volume and a substantial part of the processing workload. Infrastructure planning based only on the size of the original citizen-generated record would therefore underestimate the requirements of the complete research-data lifecycle.

This experience also favours a sample-centred and modular data model over a single flattened dataset. Different forms of enrichment have different structures, provenance and computational characteristics, but can remain related through a stable sample identity. New analytical products can consequently be added without repeatedly restructuring the complete citizen-observation dataset.

Fourth, privacy protection should be represented explicitly as a data transformation. In the case of geographical information, publishing an obfuscated coordinate without indicating that its precision has been reduced could lead downstream users to interpret the value incorrectly. Recording the applied obfuscation radius makes this transformation visible and allows users to assess whether the resulting spatial precision is appropriate for their analysis. More generally, privacy-preserving transformations should be documented as part of the provenance of the published research product.

Fifth, quality correction and privacy protection should remain separate stages. The repository should first establish the best-supported scientific representation of a record and only then derive the representation appropriate for public dissemination. This prevents privacy transformations from interfering with quality-control procedures and allows corrected data to be republished without losing the underlying provenance.

Sixth, FAIR-oriented publication is easier to sustain when it is incorporated into the normal data-processing workflow rather than added as a final project-delivery task. In ECHOREPO, canonical machine-readable resources form an intermediate boundary between the continuously evolving operational database and external publication. These resources can be regenerated, validated and deposited as new versions without exposing downstream users to changes in the internal database schema.

The same principle applies to interoperability. External discovery through infrastructures such as Zenodo and SoilWise is more robust when the repository produces stable, documented research products rather than expecting external systems to interpret its operational database directly. The publication interface therefore becomes part of the architecture rather than a secondary export function.

Finally, the deployment demonstrated the importance of designing for evolution. Citizen-science projects change during operation: data-collection procedures are refined, new error classes are discovered, analytical partners introduce new outputs, terminology and translations evolve, and external publication requirements can change. A research-data infrastructure should therefore expect schema evolution and new processing stages rather than assuming that the complete data model can be fixed at the beginning of the project.

Taken together, these observations suggest a broader design principle for citizen-science research infrastructure: preserve original context, maintain stable identifiers, separate acquisition from curation, expect progressive enrichment, make quality and privacy transformations explicit, and treat persistent publication as part of the data lifecycle. These principles are not specific to soil monitoring and provide the main basis for the architectural reusability discussed in the preceding section.